\section{The specified constant-cochain module}
\label{full:module-section}

Fix $1\leq t<n<p-1$, put
\[
 m=n-t,\qquad G=\GL_m(\Zp),\qquad U_0=P_t\times G,
\]
and use the finite field $k$ and algebraic frame rings of
Section~\ref{sec:rings}.  The coefficient induction is now complete.
We determine $H^*(P_n,B_{\mathrm{conn},G})$ together with its unit
and its action by the specified rank-$n$ constant classes.

Retain the dual complex $\cD$ of \eqref{full:dual-complex}, the
independent extension locus $Z=(N_t^m)_{\mathrm{ind}}$, and the
coefficients $R=(\cO_E/\varpi)^a$ and
$\mathscr A=(\cO^{\flat,+}/\varpi)^a$ of
Section~\ref{sec:support}.  Proposition~\ref{prop:full-frame-comparison}
identifies $\cD\otimes_kR$ with compact supports on $Z$, with the
volume line $\langle\eta\rangle$ and the shift $m+n^2$.
Choose a translation subgroup $T_h=[p]^hT$ acting trivially on
$H^*(\cD)$ as in Proposition~\ref{full:joint-completion}, and put
$d=m(t-1)$.  The derived translation action on $\cD$ is still retained.

The argument first isolates a single support and translation term by
its matrix weights.  A natural top-edge map then identifies the full
derived calculation with that term.  Connected-group duality and the
block-restriction comparison of
Proposition~\ref{prop:modification-transport} identify the surviving
constant action.  Finally, the integral classes of
Theorem~\ref{thm:stable-generators} and the coefficient unit give the
specified exterior-algebra basis.

\subsection{The parabolic weight calculation}

For $t>1$, let $\Lambda_h\simeq k[[D_1,\ldots,D_d]]$ be the
distribution algebra of $T_h$, with augmentation ideal
$\mathfrak m_{\Lambda_h}$, and put
\[
 E_h^b=\bigwedge^b
       (\mathfrak m_{\Lambda_h}/\mathfrak m_{\Lambda_h}^2)^*,
       \qquad 0\leq b\leq d.
\]
For $t=1$ this means $E_h^0=k$ and there are no other terms.
Let $\epsilon_i$ denote the character of the $i$th diagonal
Teichm\"uller entry of $G$.  The $t-1$ translation Ext generators in
row $i$ each have character $\epsilon_i$.  This follows from the
contragredient change of translation parameters; any Frobenius twists
of these characters agree on $\Fp^\times$.
Thus a weight of $E_h^b$ has the form
\begin{equation}
\label{full:exterior-weights}
 \sum_{i=1}^m e_i\epsilon_i,\qquad
 0\leq e_i\leq t-1,\qquad \sum_i e_i=b.
\end{equation}
In particular, $E_h^d$ has character $(\det)^{t-1}$.

There are two other characters in the support comparison.
For a smallest plane of dimension $r$, its compact support orientation
has weight $-t$ in each of its $r$ rows and zero in the remaining
rows, by Lemma~\ref{support:translation-quotient}.
The single volume line $\langle\eta\rangle$ has weight $+1$
in every row, by \eqref{rings:volume-character}.
Consequently the coefficient weights in a flag term are
\begin{equation}
\label{full:flag-weights}
 \lambda_i=
 \begin{cases}
  1-t+e_i,&1\leq i\leq r,\\
  1+e_i,&r<i\leq m.
 \end{cases}
\end{equation}
The affine Tate twist has no $G$-character.  The signs in
\eqref{full:flag-weights} account separately for inverse substitution
on the compact support orientation and for duality on translation
Ext; the volume has been included exactly once.

\begin{lemma}[Weight vanishing]
\label{lem:weight-vanishing}
For every $a,q,b$ one has
\begin{equation}
\label{full:weight-vanishing}
 H^a\bigl(G,H_c^q(Z;\mathscr A)\otimes
               \langle\eta\rangle\otimes E_h^b\bigr)=0
\end{equation}
unless $q=(t+2)m$ and $b=d$.
For $b=d$, extension to the ambient space induces an equivalence
\begin{equation}
\label{full:ambient-edge}
 \begin{aligned}
 &\RGamma\bigl(G,\RGamma_c(Z;\mathscr A)\otimes
                    \langle\eta\rangle\otimes E_h^d\bigr)\\
 &\qquad\xrightarrow{\ \sim\ }
 \RGamma\bigl(G,\RGamma_c(N_t^m;\mathscr A)\otimes
                    \langle\eta\rangle\otimes E_h^d\bigr).
 \end{aligned}
\end{equation}
The remaining ambient coefficient is a line with trivial $G$-action.
All assertions preserve the commuting connected-frame action.
\end{lemma}

\begin{proof}
We give the group-cohomological weight bound used in the argument.
Let $P\subset G$ be the compact stabilizer of a rational flag,
written with its successive subspaces first.  Its matrix entries
below the corresponding block diagonal are zero.
Let $P^+$ consist of its matrices whose reduction is upper
unitriangular, and let
\[
 \Delta=\{\operatorname{diag}(a_1,\ldots,a_m):
                      a_i\in\mu_{p-1}\}.
\]
The subgroup $P^+\rtimes\Delta$ has index prime to $p$ in $P$.
Indeed its index is the product of the complete flag counts in
the diagonal Levi blocks, each congruent to one modulo $p$.
Restriction to this subgroup is therefore injective on
$k$-cohomology, by restriction followed by transfer.

The pro-$p$ group $P^+$ has an ordered valued basis consisting
of elementary upper entries, principal diagonal entries, and
the allowed principal lower entries.  Give these entries the
values
\begin{equation}
\label{full:iwahori-valuation}
 \begin{array}{c|c}
  \text{entry}&\text{value}\\ \hline
  E_{ij}\ (i<j)&(j-i)/m\\
  pE_{ij}\ (i>j)\text{ in one flag block}
                    &1+(j-i)/m\\
  pE_{ii}&1 .
 \end{array}
\end{equation}
Successive $p$-powers increase the value by one.
When $m=1$ only the diagonal entry occurs.
For $m>1$ the smallest value is $1/m>1/(p-1)$.
Matrix multiplication respects the resulting filtration:
the row-index terms add along a product $E_{ij}E_{jk}$, and
the $p$-adic valuations add.  The block zero conditions are
closed under these products.  Logarithm and exponential
converge in this filtration because a term of degree $s$
has valuation at least $s/m-v_p(s!)$, tending to infinity;
the same estimates identify the required $p$-roots when their
valuation exceeds the saturation threshold.  Thus
\eqref{full:iwahori-valuation} is a saturated $p$-valuation
with values in $\frac1m\Z$.

Apply the valued-group cohomology spectral sequence of
\cite[Theorem~1.1]{Sorensen}.  In its trivial-coefficient
specialization, the $E_1$-page is the cohomology of the
finite-dimensional graded Lie algebra obtained from the
ordered basis by setting the power-shift parameter to zero.
The group-ring construction is functorial for automorphisms
preserving the valuation.  It is therefore $\Delta$-equivariant;
this equivariant use is also the construction in
\cite[\S3 and equation~(4.2)]{DLH}.
The Chevalley cochain complex gives the following sufficient
bound: every weight in $H^*(P^+,k)$ is a sum of distinct allowed
cochain-root weights
\begin{equation}
\label{full:cochain-roots}
 \epsilon_j-\epsilon_i
       \quad\text{for each allowed tangent entry }E_{ij};
\end{equation}
diagonal entries contribute zero.
This statement only bounds the weights of the abutment;
it does not assert degeneration of this spectral sequence.
It follows because its Lie cochains are exterior powers of
the dual leading-symbol space, and every subsequent page is
a subquotient.

The same bound holds with the finite coefficient
representation in \eqref{full:flag-weights}, after adding
one of its weights.  To see this without assuming a trivial
unipotent action, filter that representation by powers of the
augmentation ideal of the finite $p$-group through which
$P^+$ acts.  This is a finite, $\Delta$-stable filtration.
Its quotients have trivial $P^+$-action and decompose into
$\Delta$-characters, since $|\Delta|$ is invertible in $k$.
Their character weights occur among the original coefficient
weights.  Apply the preceding calculation to each quotient
and use the coefficient spectral sequence.
Tensoring with $R$ preserves both this filtration and its
weight vanishing.

Consider a flag with smallest plane of dimension $r<m$.
Use its finite flag resolution from
Proposition~\ref{prop:compact-support}.  At each term, Shapiro
reduces to a group $P$ as above whose first block has size $r$.
For a weight in \eqref{full:flag-weights}, the sum of the
exponents on the last $m-r$ coordinates is strictly positive.
A root internal to either of the two large blocks contributes
zero to that sum.  The allowed crossing roots are precisely
the cochain roots from an upper entry with
$i\leq r<j$; by \eqref{full:cochain-roots} each contributes
$+1$ to the last-block sum.  There are no crossing roots with
the opposite sign, since the lower crossing entries are
absent from $P$.  The resulting full weight therefore cannot
be zero as an integral weight.

This integral test also tests the actual finite torus.
Each coordinate of a sum of distinct allowed roots lies
between $-(m-1)$ and $m-1$.  Equation
\eqref{full:flag-weights} gives
\[
 |\text{each resulting exponent}|
             \leq t+m-1=n-1<p-1.
\]
The only multiple of $p-1$ in this range is zero.
Hence a trivial $\Delta$-character would require every
integral exponent to be zero, contrary to the last-block
sum.  Exactness of $\Delta$-invariants and injectivity of
restriction prove vanishing for each of these parabolic
terms.  Their finite flag resolution proves it for each
proper-plane support cohomology group.

For the ambient term put $r=m$.  Its coefficient exponents
are $1-t+e_i\leq0$, with total
\[
 \sum_i(1-t+e_i)=b-m(t-1)=b-d.
\]
Every root has total zero.  If $b<d$, the same finite-torus
bound therefore excludes a trivial weight.  If $b=d$, then
every $e_i=t-1$, and the coefficient is a line with character
\[
 (\det)^{-t}\,(\det)\,(\det)^{t-1}=1.
\]
This cancellation is an equality of $G$-characters, since
the three determinant characters are defined on $G$, not
only on $\Delta$.

These calculations prove \eqref{full:weight-vanishing}.
The map to ambient support is an isomorphism on the top
geometric cohomology group by
Proposition~\ref{prop:compact-support}; its other cohomology
groups are the proper-plane groups just killed.
The bounded hypercohomology spectral sequence of its cone
therefore proves \eqref{full:ambient-edge}.
The flag maps, weight filtrations, and torus projections
commute with the connected-frame action, which is consequently
retained throughout.
\end{proof}

\subsection{The canonical top edge}

The translation action on $H^*(\cD)$ is trivial after replacing $T$
by $T_h$.  We continue to retain its action on the complex $\cD$.
Triviality on cohomology alone does not identify that derived action
with a trivial one.

\begin{lemma}[Equivariant top-edge insertion]
\label{full:top-edge}
There is a natural $U_0$-equivariant map
\begin{equation}
\label{full:top-edge-map}
 \cD\otimes_k E_h^d[-d]\longrightarrow\RGamma(T_h,\cD).
\end{equation}
On a trivially acted-on cohomology row, it is insertion into
translation degree $d$.  The map is linear for constant stabilizer
cochains.  After applying $\RGamma(G,-)$ it is an equivalence,
as a complex with the remaining smooth $P_t$-action.
\end{lemma}

\begin{proof}
Suppose $t>1$ and write $\Lambda=\Lambda_h$.
The residue field $k$ is a perfect $\Lambda$-module, of projective
dimension $d$.  Finite Koszul duality gives
the natural tensor-Hom identification
\[
 \RHom_\Lambda(k,\cD)
   \simeq
 \RHom_\Lambda(k,\Lambda)\otimes_\Lambda^{\mathbb L}\cD.
\]
The first factor on the right has one cohomology group:
\[
 \Ext_\Lambda^i(k,\Lambda)=0\quad(i\ne d),
 \qquad
 \Ext_\Lambda^d(k,\Lambda)
       =\det(\mathfrak m_\Lambda/\mathfrak m_\Lambda^2)^*=E_h^d.
\]
Truncation therefore gives a canonical identification
$\RHom_\Lambda(k,\Lambda)\simeq E_h^d[-d]$, where $\Lambda$
acts on the line through its residue field.  Consequently
\begin{equation}
\label{full:local-duality}
 \RHom_\Lambda(k,\cD)
   \simeq
 (k\otimes_\Lambda^{\mathbb L}\cD)\otimes_k E_h^d[-d].
\end{equation}
The map $\Lambda\to k$ gives the natural map
$\cD\to k\otimes_\Lambda^{\mathbb L}\cD$.
Tensor it with $E_h^d[-d]$ and use
\eqref{full:local-duality}.  This is
\eqref{full:top-edge-map}, since rational translation cohomology
is the displayed derived Hom.

These constructions are natural for every augmented automorphism
of $\Lambda$.  In particular, the orientation identification is
the canonical top Ext identification, so it is equivariant even
when a frame change acts nonlinearly on chosen parameters.
One may use coordinates to check the row map: for a module $M$
annihilated by $\mathfrak m_\Lambda$, the Koszul differential is
zero, and the map sends $M\otimes E_h^d[-d]$ to the top Ext
summand of $\RHom_\Lambda(k,M)$.
This coordinate check defines no choices in the map itself.
Tensor-Hom duality and the augmentation are natural for all
commuting coefficient maps, including the constant-cochain action.

Filter $\cD$ by its bounded cohomology filtration.
On a row $H^q(\cD)$, the target has translation cohomology
$H^q(\cD)\otimes E_h^b$ in degrees $0\leq b\leq d$.
The row map is the just-described insertion at $b=d$.
After tensoring with $R$, Proposition~\ref{prop:full-frame-comparison}
and Lemma~\ref{lem:weight-vanishing} kill all the other columns
on applying $\RGamma(G,-)$.
The same lemma identifies the top column on source and target.
The spectral sequences are bounded, since $\cD$ is bounded,
$k$ has $\Lambda$-projective dimension $d$, and $G$ has
$p$-cohomological dimension $m^2$.
The comparison is therefore an equivalence after almost scalar
extension, hence before it by faithful flatness.
All its maps commute with $P_t$.

For $t=1$, rational representations of $T_h$ decompose into
characters, and the natural projection onto the trivial character
is exact.  It defines \eqref{full:top-edge-map} with $d=0$.
All cohomology rows are already in the trivial character,
so the projection is a quasi-isomorphism; the same argument
with $G$ and the support calculation applies.
\end{proof}

We now combine the comparisons, keeping track of their actions
and of the connected-frame limit.  Write
\[
 A_G=H^*_{\mathrm{cts}}(G,k),\qquad
 M_t^*=H^*(P_n,B_{\mathrm{conn},G}).
\]
The star in $M_t^*$ denotes its cohomological grading.
When taking its algebraic dual below, we write $\Hom_k(-,k)$
to avoid confusing these two uses.

\begin{proposition}[The surviving constant action]
\label{full:surviving-module}
After tensoring with $R$ and choosing the one-dimensional
orientation factors, there are graded identifications
\begin{equation}
\label{full:dual-poincare-module}
 \Hom_k(M_t^i,k)\otimes_kR
       \simeq A_G^{m^2-i}\otimes_kR.
\end{equation}
If a constant class on $P_n$ corresponds to a class on
$\GL_n(\Zp)$ under Proposition~\ref{prop:modification-transport},
its transpose action on the right of
\eqref{full:dual-poincare-module} is ordinary cup product by
the upper-left block restriction of that class to $G$.
\end{proposition}

\begin{proof}
Put
\[
 C_G=\RGamma\bigl(G,\RGamma(T_h,\cD)\bigr).
\]
Lemma~\ref{full:trace-descent} identifies $\RGamma(K,C_G)$ with
the dual of the coefficient complex at level $(K,G,h)$.
The proof of Proposition~\ref{full:joint-completion} gives finite
cohomology for $C_G$ before almost scalar extension, with its smooth
$P_t$-action.  Applying \eqref{full:connected-complex-limit} therefore
gives
\begin{equation}
\label{full:constant-connected-limit}
 \Hom_k\!\left(H^i(P_n,B_{\mathrm{conn},G}^{(h)}),k\right)
       \simeq H^{-i-t^2}(C_G).
\end{equation}
Here corestriction is the transpose of coefficient pullback, so its
inverse limit is exactly the dual of the connected-frame union.
The finite Mittag--Leffler calculation takes place before tensoring
with $R$.  Power transport removes the root stage while preserving
ordinary $\Fp$-cochains.

To compute $C_G$, apply Lemma~\ref{full:top-edge} and then
Proposition~\ref{prop:full-frame-comparison} after the faithful almost
extension.  The ambient map \eqref{full:ambient-edge} and
Lemma~\ref{support:translation-quotient} leave the ambient line in
geometric degree $(t+2)m$.  The three $G$-characters cancel as in
Lemma~\ref{lem:weight-vanishing}.

If $a$ is the $G$-cohomological degree, the dual cochain lies in
degree
$a+(t+2)m+d+t^2-(m+n^2)$.
It is consequently dual to primal degree
\begin{equation}
\label{full:degree-cancellation}
 i=n^2+m-(t+2)m-d-t^2-a=m^2-a.
\end{equation}
Here $n=m+t$ and $d=m(t-1)$, including $d=0$ when $t=1$.
The remaining Tate and connected orientation factors are
one-dimensional.  A choice of their bases over the fixed
geometric coefficient field gives
\eqref{full:dual-poincare-module}; changing those bases changes
the comparison by a nonzero scalar and does not change its
module assertion.

To identify the action, return to the support calculation on $Z$
before the ambient map.  Both the support comparison and the top
translation insertion preserve the actual pullback of constant
$P_n$-cochains along the middle-bundle map.  A constant class of
degree $r$ thus acts by a $P_t$-equivariant map
$C_G\to C_G[r]$ in the derived category.
The naturality for degree-shifted coefficient maps in
\eqref{full:connected-complex-limit} shows that its action in
\eqref{full:constant-connected-limit} is the underlying map on
$H^*(C_G)$ after forgetting $P_t$.
Equivalently, the connected-group spectral sequence retains only
degree $t^2$; a positive connected-filtration component would raise
that degree beyond the cohomological dimension and vanishes.
This uses no equivariant splitting of $C_G$.  Compact duality retains
the dual cup signs throughout.

After fixing the connected quotient frame,
Proposition~\ref{prop:modification-transport}, specifically
\eqref{eq:const-block-action}, identifies this ordinary action
on $[Z/G]$ with pullback from $BG$ of the actual restriction
along $g\mapsto\operatorname{diag}(g,1_t)$.
That equality comes from the slope-zero extension with its
fixed quotient and the vanishing of positive finite-coefficient
cohomology of the additive rational upper block; it is an
equality of pullback classes, so applies to their cup action
on compact supports.
The map from $Z$ to ambient support is $G$-equivariant and
hence linear for these $BG$-classes.  Thus it preserves the
now-identified action.  The middle positive bundle need only
be of standard type on $Z$; no extension of its $BP_n$-map
to dependent tuples is used.
The ambient line is $G$-trivial, so the remaining action is
ordinary cup product on $A_G$.  This proves the assertion.
\end{proof}

\subsection{The constant basis on the primal side}

Fix the classes of Theorem~\ref{thm:stable-generators} at rank $n$.
Write
\[
 \xi_i\in H^{2i-1}(P_n,k),\qquad 1\leq i\leq n,
\]
for the transported reductions of its integral matrix classes
$\alpha_{n,i}$.  They are also the leading ordinary classes of
the same continuous sphere-cochain representatives $u_i$ that
will be used in the height tests.

\begin{proposition}[The actual constant-cochain module]
\label{prop:constant-module}
For $1\leq t<n<p-1$, put $m=n-t$.  The cup map
\begin{equation}
\label{full:constant-basis}
 \bigwedge_k(e_1,\ldots,e_m)
     \longrightarrow H^*(P_n,B_{\mathrm{conn},G}),
 \qquad e_i\longmapsto \xi_i\cdot1,\qquad |e_i|=2i-1,
\end{equation}
is an isomorphism of graded $k$-algebras.
In particular, for $I\subset\{1,\ldots,m\}$ the products
\[
 \xi_I\cdot1=\left(\prod_{i\in I}\xi_i\right)\cdot1
\]
form a basis in cohomological degrees
$\sum_{i\in I}(2i-1)$ and internal degree zero.
The classes $\xi_i$ for $i>m$ act as zero.
The assertion uses the fixed rank-$n$ classes simultaneously
for every $t$.
\end{proposition}

\begin{proof}
By Theorems~\ref{thm:constant-cohomology} and
\ref{thm:stable-generators},
\[
 A_G=\bigwedge_k
       (\overline\alpha_{m,1},\ldots,\overline\alpha_{m,m}),
 \qquad \sum_{i=1}^m(2i-1)=m^2.
\]
The stable-generator theorem identifies the block restriction of
$\overline\alpha_{n,i}$ with $\overline\alpha_{m,i}$ for
$i\leq m$, and with zero for $i>m$.
Proposition~\ref{full:surviving-module} consequently identifies
the dual of the desired module, after tensoring with $R$,
with the reindexed regular module of this specified exterior
algebra.

An exterior algebra has its perfect top-degree pairing
\[
 A_G^j\otimes_k A_G^{m^2-j}
       \longrightarrow A_G^{m^2}\simeq k.
\]
Indeed the ordered monomial for a subset pairs, up to its
exterior sign, with the monomial for its complement; all other
pairings in complementary degree are zero.
Thus the reindexed dual of the regular $A_G$-module is again
its regular module, with a generator in degree zero.
This proves that the target of \eqref{full:constant-basis},
after tensoring with $R$, is free of rank one under the
displayed exterior action.

We identify its generator.  The preceding comparison shows
that $H^0(P_n,B_{\mathrm{conn},G})$ is one-dimensional over $k$:
it is finite by Proposition~\ref{full:joint-completion},
and its almost scalar extension has the length of one copy
of $R$.  The unit $1$ is nonzero in it.  More explicitly,
the algebraic coefficient rings and their frame pullbacks
are unital, so $k\cdot1$ injects into their union; this line
is $P_n$-invariant, and no negative-degree cochains can give
a boundary in degree zero.  Hence $1$ spans that
one-dimensional $k$-space.  After tensoring with $R$ it is
a basis of the degree-zero line, and therefore generates
the entire regular module.

The cup map in \eqref{full:constant-basis} is consequently an
isomorphism after the faithful almost extension.  Its kernel
and cokernel are $k$-vector spaces, and exactness and
faithfulness of that extension make both zero.
The map is a map of algebras: group cohomology with
commutative coefficients has graded-commutative cup product,
and each $\xi_i$ has odd degree, so its square vanishes
because $p$ is odd.  This proves the algebra presentation and
the asserted product basis.
The zero restrictions for $i>m$ give zero actions after
extension by Proposition~\ref{full:surviving-module}, hence
before extension by the same faithfulness.
All comparisons preserve the actual classes and powering
maps, so the proof applies to the one fixed family at rank
$n$, rather than choosing new generators at each height.
\end{proof}
